%% ============================================================================
%% Section: Evaluation
%% Included from main.tex via %% ============================================================================
%% Section: Evaluation
%% Included from main.tex via %% ============================================================================
%% Section: Evaluation
%% Included from main.tex via %% ============================================================================
%% Section: Evaluation
%% Included from main.tex via \input{sections/evaluation}
%% Owns Fig fig:latency and Fig fig:incr (referenced from \S7).
%% SKELETON: experimental design is written; all quantitative results are \TODO
%% pending measurement. See the status note.
%% ============================================================================
\section{Evaluation}\label{sec:eval}

\NOTE{Status: this section is the experimental \emph{design}. No results are
measured yet; every number is a \TODO. Do not submit without filling these. The
prose fixes the questions, baselines, datasets, and metrics so the harness has a
target; keep the ``pending'' framing until the artifact produces numbers.}

Our evaluation asks five questions. \textbf{RQ1 (query performance):} how do
Causal4j's structural and identification queries compare in latency and
throughput against alternatives? \textbf{RQ2 (index cost):} what does the
compiled index cost to build and to store, as a function of graph size and
density? \textbf{RQ3 (maintenance):} when does incremental maintenance beat full
recompilation? \textbf{RQ4 (scalability):} how does query latency grow with
$|V|$? \textbf{RQ5 (correctness):} do the answers agree with established causal
tools?

\paragraph{Baselines}
We compare four points along the ``how close to the DBMS'' axis: (a) an in-DBMS
map-of-sets adjacency with per-query DFS---the un-indexed baseline of
Table~\ref{tab:complexity}; (b) plain Cypher variable-length path queries
answering the same structural questions declaratively; (c) Neo4j GDS traversal
primitives~\cite{TODO_neo4j}; and (d) round-tripping the DAG to external causal
libraries---DoWhy~\cite{TODO_dowhy}, dagitty~\cite{TODO_dagitty},
pgmpy~\cite{TODO_pgmpy}, and causal-learn~\cite{TODO_causallearn}---which
includes export/import cost, the overhead our in-DBMS design removes.

\paragraph{Datasets}
We use standard Bayesian-network DAGs of increasing size (from small networks such
as \textsc{asia} and \textsc{sachs} through \textsc{alarm} and \textsc{hepar2} to
large ones such as \textsc{munin}), synthetic scale-free DAGs generated with a
preferential-attachment model to control $|V|$ and density independently, and the
COVID-19 running-example DAG~\cite{piccininni2020covid}. \TODO{Fix the exact
network list and the synthetic $|V|$ grid ($10^2$--$10^4$) and edge-density
settings; state the generator and seeds for reproducibility.}

\paragraph{Metrics and protocol}
For RQ1 we report median and tail latency and sustained throughput per operator
(ancestors/descendants, d-separation, adjustment-set enumeration). For RQ2, index
build time and resident memory vs.\ $|V|$ and density. For RQ3, wall-clock of an
incremental batch vs.\ a full recompile across batch sizes and edit types
(insert, delete, active/latent flip). For RQ4, latency vs.\ $|V|$ on a log--log
scale. For RQ5, exact agreement of d-separation verdicts and adjustment sets
against dagitty and causal-learn on every dataset. \TODO{State hardware, JVM
version and heap, Neo4j version, warmup and repeat counts, and report variance
(CI or IQR), per the reproducibility brief in the outline.}

\paragraph{Expected results (to be measured)}
On DAGs from \TODO{$V_{\min}$} to \TODO{$V_{\max}$} variables we expect Causal4j
to answer d-separation and adjustment-set queries \TODO{$N\times$} faster than
in-DBMS DFS and \TODO{$M\times$} faster than an external round-trip, to build its
index in \TODO{under X\,ms}, and to apply an incremental edge batch in
\TODO{Y} versus \TODO{Z} for a full recompile, while matching the reference tools
exactly on correctness (RQ5). Figures~\ref{fig:latency} and~\ref{fig:incr} report
RQ4 and RQ3 respectively.

\begin{figure}[t]\centering
  \fbox{\parbox{0.9\linewidth}{\centering \TODO{Fig.: query latency vs.\ $|V|$
  (log--log), one line per method (Causal4j index, DFS baseline, Cypher, GDS,
  external round-trip).}}}
  \caption{Query latency vs.\ graph size (RQ4). \emph{Pending measurement.}}
  \label{fig:latency}
\end{figure}

\begin{figure}[t]\centering
  \fbox{\parbox{0.9\linewidth}{\centering \TODO{Fig.: incremental update vs.\
  full recompile, by batch size and edit type; mark the crossover.}}}
  \caption{Incremental maintenance vs.\ recompilation (RQ3). \emph{Pending
  measurement.}}
  \label{fig:incr}
\end{figure}

%% Owns Fig fig:latency and Fig fig:incr (referenced from \S7).
%% SKELETON: experimental design is written; all quantitative results are \TODO
%% pending measurement. See the status note.
%% ============================================================================
\section{Evaluation}\label{sec:eval}

\NOTE{Status: this section is the experimental \emph{design}. No results are
measured yet; every number is a \TODO. Do not submit without filling these. The
prose fixes the questions, baselines, datasets, and metrics so the harness has a
target; keep the ``pending'' framing until the artifact produces numbers.}

Our evaluation asks five questions. \textbf{RQ1 (query performance):} how do
Causal4j's structural and identification queries compare in latency and
throughput against alternatives? \textbf{RQ2 (index cost):} what does the
compiled index cost to build and to store, as a function of graph size and
density? \textbf{RQ3 (maintenance):} when does incremental maintenance beat full
recompilation? \textbf{RQ4 (scalability):} how does query latency grow with
$|V|$? \textbf{RQ5 (correctness):} do the answers agree with established causal
tools?

\paragraph{Baselines}
We compare four points along the ``how close to the DBMS'' axis: (a) an in-DBMS
map-of-sets adjacency with per-query DFS---the un-indexed baseline of
Table~\ref{tab:complexity}; (b) plain Cypher variable-length path queries
answering the same structural questions declaratively; (c) Neo4j GDS traversal
primitives~\cite{TODO_neo4j}; and (d) round-tripping the DAG to external causal
libraries---DoWhy~\cite{TODO_dowhy}, dagitty~\cite{TODO_dagitty},
pgmpy~\cite{TODO_pgmpy}, and causal-learn~\cite{TODO_causallearn}---which
includes export/import cost, the overhead our in-DBMS design removes.

\paragraph{Datasets}
We use standard Bayesian-network DAGs of increasing size (from small networks such
as \textsc{asia} and \textsc{sachs} through \textsc{alarm} and \textsc{hepar2} to
large ones such as \textsc{munin}), synthetic scale-free DAGs generated with a
preferential-attachment model to control $|V|$ and density independently, and the
COVID-19 running-example DAG~\cite{piccininni2020covid}. \TODO{Fix the exact
network list and the synthetic $|V|$ grid ($10^2$--$10^4$) and edge-density
settings; state the generator and seeds for reproducibility.}

\paragraph{Metrics and protocol}
For RQ1 we report median and tail latency and sustained throughput per operator
(ancestors/descendants, d-separation, adjustment-set enumeration). For RQ2, index
build time and resident memory vs.\ $|V|$ and density. For RQ3, wall-clock of an
incremental batch vs.\ a full recompile across batch sizes and edit types
(insert, delete, active/latent flip). For RQ4, latency vs.\ $|V|$ on a log--log
scale. For RQ5, exact agreement of d-separation verdicts and adjustment sets
against dagitty and causal-learn on every dataset. \TODO{State hardware, JVM
version and heap, Neo4j version, warmup and repeat counts, and report variance
(CI or IQR), per the reproducibility brief in the outline.}

\paragraph{Expected results (to be measured)}
On DAGs from \TODO{$V_{\min}$} to \TODO{$V_{\max}$} variables we expect Causal4j
to answer d-separation and adjustment-set queries \TODO{$N\times$} faster than
in-DBMS DFS and \TODO{$M\times$} faster than an external round-trip, to build its
index in \TODO{under X\,ms}, and to apply an incremental edge batch in
\TODO{Y} versus \TODO{Z} for a full recompile, while matching the reference tools
exactly on correctness (RQ5). Figures~\ref{fig:latency} and~\ref{fig:incr} report
RQ4 and RQ3 respectively.

\begin{figure}[t]\centering
  \fbox{\parbox{0.9\linewidth}{\centering \TODO{Fig.: query latency vs.\ $|V|$
  (log--log), one line per method (Causal4j index, DFS baseline, Cypher, GDS,
  external round-trip).}}}
  \caption{Query latency vs.\ graph size (RQ4). \emph{Pending measurement.}}
  \label{fig:latency}
\end{figure}

\begin{figure}[t]\centering
  \fbox{\parbox{0.9\linewidth}{\centering \TODO{Fig.: incremental update vs.\
  full recompile, by batch size and edit type; mark the crossover.}}}
  \caption{Incremental maintenance vs.\ recompilation (RQ3). \emph{Pending
  measurement.}}
  \label{fig:incr}
\end{figure}

%% Owns Fig fig:latency and Fig fig:incr (referenced from \S7).
%% SKELETON: experimental design is written; all quantitative results are \TODO
%% pending measurement. See the status note.
%% ============================================================================
\section{Evaluation}\label{sec:eval}

\NOTE{Status: this section is the experimental \emph{design}. No results are
measured yet; every number is a \TODO. Do not submit without filling these. The
prose fixes the questions, baselines, datasets, and metrics so the harness has a
target; keep the ``pending'' framing until the artifact produces numbers.}

Our evaluation asks five questions. \textbf{RQ1 (query performance):} how do
Causal4j's structural and identification queries compare in latency and
throughput against alternatives? \textbf{RQ2 (index cost):} what does the
compiled index cost to build and to store, as a function of graph size and
density? \textbf{RQ3 (maintenance):} when does incremental maintenance beat full
recompilation? \textbf{RQ4 (scalability):} how does query latency grow with
$|V|$? \textbf{RQ5 (correctness):} do the answers agree with established causal
tools?

\paragraph{Baselines}
We compare four points along the ``how close to the DBMS'' axis: (a) an in-DBMS
map-of-sets adjacency with per-query DFS---the un-indexed baseline of
Table~\ref{tab:complexity}; (b) plain Cypher variable-length path queries
answering the same structural questions declaratively; (c) Neo4j GDS traversal
primitives~\cite{TODO_neo4j}; and (d) round-tripping the DAG to external causal
libraries---DoWhy~\cite{TODO_dowhy}, dagitty~\cite{TODO_dagitty},
pgmpy~\cite{TODO_pgmpy}, and causal-learn~\cite{TODO_causallearn}---which
includes export/import cost, the overhead our in-DBMS design removes.

\paragraph{Datasets}
We use standard Bayesian-network DAGs of increasing size (from small networks such
as \textsc{asia} and \textsc{sachs} through \textsc{alarm} and \textsc{hepar2} to
large ones such as \textsc{munin}), synthetic scale-free DAGs generated with a
preferential-attachment model to control $|V|$ and density independently, and the
COVID-19 running-example DAG~\cite{piccininni2020covid}. \TODO{Fix the exact
network list and the synthetic $|V|$ grid ($10^2$--$10^4$) and edge-density
settings; state the generator and seeds for reproducibility.}

\paragraph{Metrics and protocol}
For RQ1 we report median and tail latency and sustained throughput per operator
(ancestors/descendants, d-separation, adjustment-set enumeration). For RQ2, index
build time and resident memory vs.\ $|V|$ and density. For RQ3, wall-clock of an
incremental batch vs.\ a full recompile across batch sizes and edit types
(insert, delete, active/latent flip). For RQ4, latency vs.\ $|V|$ on a log--log
scale. For RQ5, exact agreement of d-separation verdicts and adjustment sets
against dagitty and causal-learn on every dataset. \TODO{State hardware, JVM
version and heap, Neo4j version, warmup and repeat counts, and report variance
(CI or IQR), per the reproducibility brief in the outline.}

\paragraph{Expected results (to be measured)}
On DAGs from \TODO{$V_{\min}$} to \TODO{$V_{\max}$} variables we expect Causal4j
to answer d-separation and adjustment-set queries \TODO{$N\times$} faster than
in-DBMS DFS and \TODO{$M\times$} faster than an external round-trip, to build its
index in \TODO{under X\,ms}, and to apply an incremental edge batch in
\TODO{Y} versus \TODO{Z} for a full recompile, while matching the reference tools
exactly on correctness (RQ5). Figures~\ref{fig:latency} and~\ref{fig:incr} report
RQ4 and RQ3 respectively.

\begin{figure}[t]\centering
  \fbox{\parbox{0.9\linewidth}{\centering \TODO{Fig.: query latency vs.\ $|V|$
  (log--log), one line per method (Causal4j index, DFS baseline, Cypher, GDS,
  external round-trip).}}}
  \caption{Query latency vs.\ graph size (RQ4). \emph{Pending measurement.}}
  \label{fig:latency}
\end{figure}

\begin{figure}[t]\centering
  \fbox{\parbox{0.9\linewidth}{\centering \TODO{Fig.: incremental update vs.\
  full recompile, by batch size and edit type; mark the crossover.}}}
  \caption{Incremental maintenance vs.\ recompilation (RQ3). \emph{Pending
  measurement.}}
  \label{fig:incr}
\end{figure}

%% Owns Fig fig:latency and Fig fig:incr (referenced from \S7).
%% SKELETON: experimental design is written; all quantitative results are \TODO
%% pending measurement. See the status note.
%% ============================================================================
\section{Evaluation}\label{sec:eval}

\NOTE{Status: this section is the experimental \emph{design}. No results are
measured yet; every number is a \TODO. Do not submit without filling these. The
prose fixes the questions, baselines, datasets, and metrics so the harness has a
target; keep the ``pending'' framing until the artifact produces numbers.}

Our evaluation asks five questions. \textbf{RQ1 (query performance):} how do
Causal4j's structural and identification queries compare in latency and
throughput against alternatives? \textbf{RQ2 (index cost):} what does the
compiled index cost to build and to store, as a function of graph size and
density? \textbf{RQ3 (maintenance):} when does incremental maintenance beat full
recompilation? \textbf{RQ4 (scalability):} how does query latency grow with
$|V|$? \textbf{RQ5 (correctness):} do the answers agree with established causal
tools?

\paragraph{Baselines}
We compare four points along the ``how close to the DBMS'' axis: (a) an in-DBMS
map-of-sets adjacency with per-query DFS---the un-indexed baseline of
Table~\ref{tab:complexity}; (b) plain Cypher variable-length path queries
answering the same structural questions declaratively; (c) Neo4j GDS traversal
primitives~\cite{TODO_neo4j}; and (d) round-tripping the DAG to external causal
libraries---DoWhy~\cite{TODO_dowhy}, dagitty~\cite{TODO_dagitty},
pgmpy~\cite{TODO_pgmpy}, and causal-learn~\cite{TODO_causallearn}---which
includes export/import cost, the overhead our in-DBMS design removes.

\paragraph{Datasets}
We use standard Bayesian-network DAGs of increasing size (from small networks such
as \textsc{asia} and \textsc{sachs} through \textsc{alarm} and \textsc{hepar2} to
large ones such as \textsc{munin}), synthetic scale-free DAGs generated with a
preferential-attachment model to control $|V|$ and density independently, and the
COVID-19 running-example DAG~\cite{piccininni2020covid}. \TODO{Fix the exact
network list and the synthetic $|V|$ grid ($10^2$--$10^4$) and edge-density
settings; state the generator and seeds for reproducibility.}

\paragraph{Metrics and protocol}
For RQ1 we report median and tail latency and sustained throughput per operator
(ancestors/descendants, d-separation, adjustment-set enumeration). For RQ2, index
build time and resident memory vs.\ $|V|$ and density. For RQ3, wall-clock of an
incremental batch vs.\ a full recompile across batch sizes and edit types
(insert, delete, active/latent flip). For RQ4, latency vs.\ $|V|$ on a log--log
scale. For RQ5, exact agreement of d-separation verdicts and adjustment sets
against dagitty and causal-learn on every dataset. \TODO{State hardware, JVM
version and heap, Neo4j version, warmup and repeat counts, and report variance
(CI or IQR), per the reproducibility brief in the outline.}

\paragraph{Expected results (to be measured)}
On DAGs from \TODO{$V_{\min}$} to \TODO{$V_{\max}$} variables we expect Causal4j
to answer d-separation and adjustment-set queries \TODO{$N\times$} faster than
in-DBMS DFS and \TODO{$M\times$} faster than an external round-trip, to build its
index in \TODO{under X\,ms}, and to apply an incremental edge batch in
\TODO{Y} versus \TODO{Z} for a full recompile, while matching the reference tools
exactly on correctness (RQ5). Figures~\ref{fig:latency} and~\ref{fig:incr} report
RQ4 and RQ3 respectively.

\begin{figure}[t]\centering
  \fbox{\parbox{0.9\linewidth}{\centering \TODO{Fig.: query latency vs.\ $|V|$
  (log--log), one line per method (Causal4j index, DFS baseline, Cypher, GDS,
  external round-trip).}}}
  \caption{Query latency vs.\ graph size (RQ4). \emph{Pending measurement.}}
  \label{fig:latency}
\end{figure}

\begin{figure}[t]\centering
  \fbox{\parbox{0.9\linewidth}{\centering \TODO{Fig.: incremental update vs.\
  full recompile, by batch size and edit type; mark the crossover.}}}
  \caption{Incremental maintenance vs.\ recompilation (RQ3). \emph{Pending
  measurement.}}
  \label{fig:incr}
\end{figure}
